\section*{Introdução}

O documento a desenvolver deve cumprir os seguintes requisitos:

\begin{itemize}

\item O documento \textbf{não pode ter mais de 15 páginas}, incluindo tabelas, figuras, referências, equações ou anexos. Caso se exceda o limite de páginas, a nota será penalizada.

\item O documento deve conter as seguintes secções:

\begin{itemize}

\item \textbf{Resumo}: não pode ultrapassar as 250 palavras e deve conter os aspetos mais relevantes que os autores pretendem destacar sobre o trabalho. Por exemplo, os autores podem utilizar o resumo para indicar brevemente os procedimentos utilizados e os resultados mais relevantes, tais como o erro relativo ao comparar os resultados obtidos por diferentes métodos.

\item \textbf{Introdução}: na secção de introdução deve identificar-se a aeronave escolhida, contextualizar a aeronave e a sua utilização, e identificar os perfis aerodinâmicos que constituem a asa. Caso a asa disponha de diferentes perfis ao longo da envergadura, utilizar-se-á o perfil correspondente à raiz da asa como caso de referência.

\item \textbf{Procedimentos}: nesta secção devem explicar-se os procedimentos ou métodos seguidos para o estudo aerodinâmico do perfil que constitui a asa (método dos perfis delgados, método dos painéis, simulação CFD). Informação adicional sobre esta secção é apresentada no ponto correspondente.

\item \textbf{Resultados}: devem apresentar-se e explicar-se os resultados obtidos através dos três métodos. Os autores devem demonstrar que compreendem as vantagens e desvantagens de cada método, bem como o seu alcance. Os resultados solicitados são especificados no ponto correspondente.

\item \textbf{Conclusões}: nesta secção os estudantes devem expor as conclusões que consideram mais relevantes do trabalho.

\end{itemize}

\end{itemize}
